\section{Objective and outcome}
Synthetic closed-orbit comparison for typesetting only; no solver was run.
\section{Physical basis and model}
For a linear transverse map, the symplectic condition is
\begin{equation}
 M^{\mathsf T}JM=J,\qquad
 J=\begin{pmatrix}0&1\\-1&0\end{pmatrix}.
\label{eq:map}
\end{equation}
\begin{table}[htbp]
\centering
\begin{tabularx}{\linewidth}{Y r r Y}
\toprule Quantity & Model & Example COF & Qualification \\
\midrule
Full horizontal tune [turns/turn] & 0.800 & 0.800 & Example values, no acceptance claim \\
Full vertical tune [turns/turn] & 1.700 & 1.700 & Same coordinate convention \\
Horizontal closure residual [m] & 0 & $10^{-10}$ & Illustrative numerical residual only \\
\bottomrule
\end{tabularx}
\caption{Example comparison with long labels and explicit units.}\label{tab:comparison}
\end{table}
\begin{table}[htbp]
\centering
\begin{tabularx}{\linewidth}{l Y Y}
\toprule Attempt & Status & Interpretation \\
\midrule
001 & Failed: incomplete evidence & Retained for provenance; no accepted tune \\
002 & Completed illustration & Does not establish a physical result \\
\bottomrule
\end{tabularx}
\caption{Example attempt history.}
\end{table}
\begin{figure}[htbp]
\centering
\includegraphics[width=\linewidth,height=0.20\textheight,keepaspectratio]{figure.png}
\caption{Existing image checking caption and figure layout only.}
\end{figure}
\section{Conclusion and limitations}
Table~\ref{tab:comparison} illustrates comparison layout. Equation~\ref{eq:map}
uses the linear-map convention; no convergence or stability result is claimed.
